%! Scatterplots and Correlation — Unit 2, Lesson 2.2 — Homework
% GRADUAL RELEASE: this is the lesson's SCORED practice and the LAST component in the packet.
% Paper homework is the DEFAULT; DeltaMath is an occasional override that replaces this page and
% strikes its score cell on the cover. Started in the last eight minutes of the period.
% ONE PAGE (compact budget), five items in two parts: A1–A2 on the plot, B1–B3 on r.
% THIRD CONTEXT: the notes teach on cars, Guided Practice runs on a city pool, AP Practice on
% screen time and sleep; Part A is a basketball team — fresh data again.
%   practice (h/week): 2 3 3 4 5 5 6 7 7 8 9 10 11 12
%   free-throw %     : 52 55 60 58 63 57 66 64 70 68 74 72 50 80      r = 0.60
%   (the 11-hour, 50% player is the unusual point; without them r is about 0.95)
% B2 is the target misconception (r near 0 read as "no relationship"); B3 is causation.
% \answerspace{H}{} reserves exactly H in the blank; the key passes the answer as the second
% argument, so the two files paginate identically by construction.
\documentclass[12pt]{article}
\usepackage{apstats-article}
\usepackage{apstats-boxes}
\newcommand{\answerspace}[2]{\par\nopagebreak\noindent\begin{minipage}[t][#1][t]{\linewidth}%
  \color{keyred}\bfseries #2\end{minipage}\par}

\begin{document}

\pageheader{Unit 2, Lesson 2.2}{Homework}
% NAMESTRIP: no \namedateperiod here — the name row lives on the cover only.

\vspace{0.04in}

% Authored BYTE-IDENTICALLY in homework/ and homework_key/.
\boxguard[8]
\begin{remindbox}
\small
\textbf{This is your graded homework.} Scored; due the first class after two study halls.
\end{remindbox}

\vspace{0.04in}

\boxguard[20]
\begin{notesbox}{Part A --- Practice and Free Throws}
\begin{minipage}[c]{0.44\linewidth}\raggedright
A coach recorded each of 14 players' hours of free-throw practice per week and their
free-throw percentage for the season. Technology gives $r = 0.60$.
\end{minipage}\hfill
\begin{minipage}[c]{0.54\linewidth}\centering
\begin{tikzpicture}[x=0.52cm, y=0.05cm]
  \foreach \y in {50,60,70,80} {
    \draw[linegray, very thin] (0,\y) -- (13,\y);
    \node[left, font=\tiny] at (0,\y) {\y};}
  \draw[thick] (0,45) -- (13.2,45);
  \draw[thick] (0,45) -- (0,84);
  \foreach \x in {0,2,4,6,8,10,12} {
    \draw (\x,44.3) -- (\x,45.7); \node[below, font=\tiny] at (\x,44.5) {\x};}
  \foreach \h/\p in {2/52, 3/55, 3/60, 4/58, 5/63, 5/57, 6/66, 7/64, 7/70, 8/68, 9/74,
                     10/72, 11/50, 12/80}
    \fill[navy] (\h,\p) circle (1.5pt);
  \node[font=\tiny] at (6.5,33.5) {Practice (hours per week)};
  \node[rotate=90, font=\tiny] at (-1.3,64.5) {Free throws (\%)};
\end{tikzpicture}
\end{minipage}
\begin{enumerate}[leftmargin=1.9em, itemsep=2pt, topsep=4pt]
  \item Explanatory variable: \blank{4.2cm} \quad Response: \blank{4.2cm}
  \item Describe the association in context (direction, unusual point, form, strength).
        \answerspace{1.2cm}{}
\end{enumerate}
\end{notesbox}

\vspace{0.04in}

\boxguard[14]
\begin{notesbox}{Part B --- What $r$ Can and Cannot Tell You}
\begin{enumerate}[leftmargin=1.9em, itemsep=2pt, topsep=2pt]
  \item \textbf{AP-style multiple choice.} Which correlation shows the \emph{strongest}
        linear association? Circle one, then justify it in a sentence.\par
        \textbf{(A)} $r = 0.45$ \quad \textbf{(B)} $r = -0.91$ \quad \textbf{(C)} $r = 0$
        \quad \textbf{(D)} $r = 0.78$ \quad \textbf{(E)} $r = -0.20$
        \answerspace{0.7cm}{}
  \item For 40 people aged 5 to 85, age vs.\ reaction time makes a clear U (young adults
        fastest), and $r = 0.05$. Jordan says age and reaction time are unrelated. Is he right?
        \answerspace{1.1cm}{}
  \item Across 50 towns, $r = 0.92$ between the number of libraries and the number of crimes.
        Should a town close libraries to cut crime? Explain.
        \answerspace{1.1cm}{}
\end{enumerate}
\end{notesbox}

\vspace{0.04in}

% The next-lesson preview closes the packet.
\boxguard[4]
\begin{spiralbox}
\small
\textbf{Next lesson: Regression Lines and Residuals.} A linear pattern earns a line: we use it to
\emph{predict} the response, then measure how far each point misses.
\end{spiralbox}

\end{document}
